\documentclass[11pt, a4paper]{article}
\usepackage[a4paper, margin=2.5cm]{geometry}
\usepackage{fontspec}
\usepackage[indonesian, bidi=basic, provide=*]{babel}
\babelprovide[import, onchar=ids fonts]{english}
\babelfont{rm}{Noto Sans}
\babelfont{tt}{Noto Sans Mono}
\usepackage{enumitem}
\setlist[itemize]{label=-}
\usepackage{amsmath, amssymb}
\usepackage{xcolor}
\usepackage{hyperref}

\hypersetup{
    colorlinks=true,
    linkcolor=blue,
    urlcolor=blue
}

\definecolor{emtcomment}{RGB}{0, 128, 0}
\definecolor{emtcommand}{RGB}{180, 0, 0}

\newcommand{\emtcmd}[1]{\noindent\texttt{\textcolor{emtcommand}{>}\ #1}\par\smallskip}
\newenvironment{emtoutput}{\begin{verbatim}}{\end{verbatim}}

\begin{document}

\section*{EMT Kalkulus}

\subsection*{Kalkulus dengan EMT}

Materi Kalkulus mencakup di antaranya:
\begin{itemize}
    \item Fungsi (fungsi aljabar, trigonometri, eksponensial, logaritma, komposisi fungsi),
    \item Limit Fungsi,
    \item Turunan Fungsi,
    \item Integral Tak Tentu,
    \item Integral Tentu dan Aplikasinya,
    \item Barisan dan Deret (kekonvergenan barisan dan deret).
\end{itemize}

\textbf{Catatan Sintaks EMT:} Kata kunci inti program Euler Math Toolbox (seperti \texttt{function}, \texttt{plot2d}, \texttt{limit}, \texttt{diff}, \texttt{integrate}, \texttt{iterate}, \texttt{sequence}, \texttt{sum}) merupakan perintah baku perangkat lunak, sedangkan seluruh penjelasan, komentar kode, nama fungsi buatan, dan penyelesaian latihan berikut telah disesuaikan ke dalam Bahasa Indonesia.

Terdapat beberapa cara mendefinisikan fungsi pada EMT, yakni:
\begin{itemize}
    \item Menggunakan format \texttt{nama\_fungsi := rumus\_fungsi} (untuk fungsi numerik),
    \item Menggunakan format \texttt{nama\_fungsi \&= rumus\_fungsi} (untuk fungsi simbolik, dapat dihitung numerik),
    \item Menggunakan format \texttt{nama\_fungsi \&\&= rumus\_fungsi} (untuk fungsi simbolik murni),
    \item Fungsi sebagai program EMT.
\end{itemize}
Setiap format harus diawali dengan perintah \texttt{function}.

\emtcmd{function f(x) := 2*x\^{}2+exp(sin(x)) \textcolor{emtcomment}{// fungsi numerik}}
\emtcmd{f(0), f(1), f(pi)}
\begin{verbatim}
    1
    4.31977682472
    20.7392088022
\end{verbatim}

\emtcmd{function g(x) := sqrt(x\^{}2-3*x)/(x+1)}
\emtcmd{g(3), g(0)}
\begin{verbatim}
    0
    0
\end{verbatim}

Jika \texttt{g(1)} dijalankan akan menghasilkan error karena akar bilangan negatif (imajiner).

\emtcmd{f(g(5)) \textcolor{emtcomment}{// komposisi fungsi}}
\begin{verbatim}
    2.20920171961
\end{verbatim}

\emtcmd{g(f(5))}
\begin{verbatim}
    0.950898070639
\end{verbatim}

\emtcmd{function h(x) := f(g(x)) \textcolor{emtcomment}{// definisi komposisi fungsi}}
\emtcmd{h(5) \textcolor{emtcomment}{// sama dengan f(g(5))}}
\begin{verbatim}
    2.20920171961
\end{verbatim}

\emtcmd{f(0:10) \textcolor{emtcomment}{// menghitung nilai f(0), f(1), ..., f(10) sekaligus}}
\begin{verbatim}
    [1,  4.31978,  10.4826,  19.1516,  32.4692,  50.3833,  72.7562,
    99.929,  130.69,  163.51,  200.58]
\end{verbatim}

\emtcmd{fmap(0:10) \textcolor{emtcomment}{// berlaku untuk semua fungsi agar menerima input vektor}}
\begin{verbatim}
    [1,  4.31978,  10.4826,  19.1516,  32.4692,  50.3833,  72.7562,
    99.929,  130.69,  163.51,  200.58]
\end{verbatim}

Mendefinisikan fungsi sepotong-sepotong (piecewise): $f(x) = x^3$ untuk $x > 0$, dan $x^2$ untuk $x \le 0$. Kata kunci \texttt{"map"} digunakan agar fungsi menerima vektor sebagai input.
\begin{verbatim}
>function map f(x)
  if x>0 then return x^3
  else return x^2
  endif;
endfunction
\end{verbatim}

\emtcmd{f(1), f(-2), f(-5:5)}
\begin{verbatim}
    1
    4
    [25,  16,  9,  4,  1,  0,  1,  8,  27,  64,  125]
\end{verbatim}

\emtcmd{aspect(1.5); plot2d("f(x)",-5,5):}

\emtcmd{function f(x) \&= 2*E\^{}x \textcolor{emtcomment}{// fungsi simbolik}}
$$2 e^x$$

\emtcmd{function g(x) \&= 3*x+1}
$$3x + 1$$

\emtcmd{function h(x) \&= f(g(x)) \textcolor{emtcomment}{// komposisi fungsi simbolik}}
$$2 e^{3x + 1}$$

\subsubsection*{1. Fungsi Polinomial (Aljabar): $f_1(x) = x^3 - 4x + 1$}
\emtcmd{function f1(x) := x\^{}3 - 4*x + 1}
\emtcmd{f1(2), f1(-2:2)}
\begin{verbatim}
    1
    [1,  4,  1,  -2,  1]
\end{verbatim}
\emtcmd{plot2d("f1(x)", -3, 3, title="1. Fungsi Polinomial: x\^{}3 - 4x + 1"):}

\subsubsection*{2. Fungsi Rasional: $f_2(x) = \frac{2x + 1}{x^2 + 1}$}
\emtcmd{function f2(x) := (2*x + 1)/(x\^{}2 + 1)}
\emtcmd{f2(0), f2(-3:3)}
\begin{verbatim}
    1
    [-0.5,  -0.6,  -0.5,  1,  1.5,  1,  0.7]
\end{verbatim}
\emtcmd{plot2d("f2(x)", -5, 5, title="2. Fungsi Rasional: (2x+1)/(x\^{}2+1)"):}

\subsubsection*{3. Fungsi Trigonometri: $f_3(x) = 2\sin(x) + \cos(2x)$}
\emtcmd{function f3(x) := 2*sin(x) + cos(2*x)}
\emtcmd{f3(pi/2), f3(0:pi/2:2*pi)}
\begin{verbatim}
    1
    [1,  1,  1,  -3,  1]
\end{verbatim}
\emtcmd{plot2d("f3(x)", -2*pi, 2*pi, title="3. Fungsi Trigonometri: 2sin(x) + cos(2x)"):}

\subsubsection*{4. Fungsi Eksponensial: $f_4(x) = 3e^{-0.5x}$}
\emtcmd{function f4(x) := 3*exp(-0.5*x)}
\emtcmd{f4(0), f4(0:5)}
\begin{verbatim}
    3
    [3,  1.81959,  1.10364,  0.66939,  0.406006,  0.246255]
\end{verbatim}
\emtcmd{plot2d("f4(x)", -2, 6, title="4. Fungsi Eksponensial: 3e\^{}(-0.5x)"):}

\subsubsection*{5. Fungsi Logaritma: $f_5(x) = \log(x^2 + 1)$}
\emtcmd{function f5(x) := log(x\^{}2 + 1)}
\emtcmd{f5(0), f5(1:5)}
\begin{verbatim}
    0
    [0.693147,  1.60944,  2.30259,  2.83321,  3.2581]
\end{verbatim}
\emtcmd{plot2d("f5(x)", -4, 4, title="5. Fungsi Logaritma: ln(x\^{}2 + 1)"):}

\subsubsection*{6. Fungsi Dua Variabel: $f_6(x, y) = \sin(x)\cos(y)$}
\emtcmd{function f6(x, y) := sin(x)*cos(y)}
\emtcmd{f6(pi/2, 0) \textcolor{emtcomment}{// nilai tunggal}}
\begin{verbatim}
    1
\end{verbatim}
\emtcmd{x\_vec = [0, pi/4, pi/2]; y\_vec = [0, pi/4, pi/2]; f6(x\_vec, y\_vec) \textcolor{emtcomment}{// nilai vektor}}
\begin{verbatim}
    Variable x not found!
    Error in:
    x_vec = [0, pi/4, pi/2]; y_vec = [0, pi/4, pi/2]; f6(x_vec, y_ ...
          ^
\end{verbatim}
\emtcmd{plot3d("f6(x,y)", xmin=-pi, xmax=pi, ymin=-pi, ymax=pi, title="Fungsi Dua Variabel: sin(x)cos(y)"):}

\subsection*{Menghitung Limit}

Perhitungan limit pada EMT menggunakan fungsi Maxima \texttt{"limit"}. Nilai limit dapat dihitung di titik tertentu, tak hingga (\texttt{inf}, \texttt{minf}), maupun limit sepihak (\texttt{plus}, \texttt{minus}).

\emtcmd{\$showev('limit((x\^{}3-13*x\^{}2+51*x-63)/(x\^{}3-4*x\^{}2-3*x+18),x,3))}

Fungsi diskontinu di $x = 3$. Berikut grafiknya:
\emtcmd{aspect(1.5); plot2d("(x\^{}3-13*x\^{}2+51*x-63)/(x\^{}3-4*x\^{}2-3*x+18)",0,4); plot2d(3,-4/5,>points,style="ow",>add):}

\emtcmd{\$showev('limit(2*x*sin(x)/(1-cos(x)),x,0))}
\emtcmd{plot2d("2*x*sin(x)/(1-cos(x))",-pi,pi); plot2d(0,4,>points,style="ow",>add):}

\emtcmd{\$showev('limit(cot(7*h)/cot(5*h),h,0))}
\emtcmd{plot2d("cot(7*x)/cot(5*x)",-0.001,0.001); plot2d(0,5/7,>points,style="ow",>add):}

Melengkapi grafik untuk contoh-contoh limit pada modul:

\emtcmd{\$showev('limit(((x/8)\^{}(1/3)-1)/(x-8),x,8))}
\emtcmd{plot2d("((x/8)\^{}(1/3)-1)/(x-8)",0,16); plot2d(8,1/24,>points,style="ow",>add):}

\emtcmd{\$showev('limit((x\^{}2-3*x-10)/(x-5),x,5))}
\emtcmd{plot2d("(x\^{}2-3*x-10)/(x-5)",0,8); plot2d(5,7,>points,style="ow",>add):}

\emtcmd{\$showev('limit(abs(x-1)/(x-1),x,1,minus))}
\emtcmd{plot2d("abs(x-1)/(x-1)",-1,3); plot2d(1,-1,>points,style="ow",>add):}

\emtcmd{\$showev('limit(sin(x\^{}3)/x,x,0))}
\emtcmd{plot2d("sin(x\^{}3)/x",-2,2); plot2d(0,0,>points,style="ow",>add):}

\emtcmd{\$showev('limit((x\^{}2-9)/(2*x\^{}2-5*x-3),x,3))}
\emtcmd{plot2d("(x\^{}2-9)/(2*x\^{}2-5*x-3)",0,5); plot2d(3,6/7,>points,style="ow",>add):}

\emtcmd{\$showev('limit((1-cos(x))/x,x,0))}
\emtcmd{plot2d("(1-cos(x))/x",-pi,pi); plot2d(0,0,>points,style="ow",>add):}

\emtcmd{\$showev('limit((x\^{}2+abs(x))/(x\^{}2-abs(x)),x,0))}
\emtcmd{plot2d("(x\^{}2+abs(x))/(x\^{}2-abs(x))",-2,2); plot2d(0,-1,>points,style="ow",>add):}

\emtcmd{\$showev('limit((1+x)\^{}(1/x),x,0))}
\emtcmd{plot2d("(1+x)\^{}(1/x)",-0.5,2); plot2d(0,exp(1),>points,style="ow",>add):}

\emtcmd{\$showev('limit((E\^{}x-E\^{}2)/(x-2),x,2))}
\emtcmd{plot2d("(exp(x)-exp(2))/(x-2)",0,4); plot2d(2,exp(2),>points,style="ow",>add):}

\subsubsection*{Latihan 2: Limit 5 Jenis Fungsi di Titik Tertentu dan Tak Hingga}

\textbf{1. Fungsi Rasional:} $g_1(x) = \frac{3x^2 - 2x + 1}{x^2 + 4}$
\emtcmd{function g1(x) \&= (3*x\^{}2 - 2*x + 1)/(x\^{}2 + 4)}
\emtcmd{\$showev('limit(g1(x), x, 2))}
\emtcmd{\$showev('limit(g1(x), x, inf))}
\emtcmd{plot2d("g1(x)", -10, 20, title="Limit g1(x) -> 3 saat x -> inf"):}

\textbf{2. Fungsi Irasional (Akar):} $g_2(x) = \frac{\sqrt{x + 4} - 2}{x}$
\emtcmd{function g2(x) \&= (sqrt(x + 4) - 2)/x}
\emtcmd{\$showev('limit(g2(x), x, 0))}
\emtcmd{\$showev('limit(g2(x), x, inf))}
\emtcmd{plot2d("g2(x)", -3, 15); plot2d(0, 1/4, >points, style="ow", >add):}

\textbf{3. Fungsi Trigonometri:} $g_3(x) = \frac{\tan(3x)}{\sin(2x)}$
\emtcmd{function g3(x) \&= tan(3*x)/sin(2*x)}
\emtcmd{\$showev('limit(g3(x), x, 0))}
\emtcmd{plot2d("g3(x)", -0.4, 0.4); plot2d(0, 3/2, >points, style="ow", >add):}

\textbf{4. Fungsi Eksponensial:} $g_4(x) = \frac{e^{2x} - 1}{x}$
\emtcmd{function g4(x) \&= (exp(2*x) - 1)/x}
\emtcmd{\$showev('limit(g4(x), x, 0))}
\emtcmd{\$showev('limit(g4(x), x, minf))}
\emtcmd{plot2d("g4(x)", -3, 2); plot2d(0, 2, >points, style="ow", >add):}

\textbf{5. Fungsi Logaritma:} $g_5(x) = \frac{\log(1 + 3x)}{x}$
\emtcmd{function g5(x) \&= log(1 + 3*x)/x}
\emtcmd{\$showev('limit(g5(x), x, 0))}
\emtcmd{\$showev('limit(g5(x), x, inf))}
\emtcmd{plot2d("g5(x)", -0.2, 10); plot2d(0, 3, >points, style="ow", >add):}

\subsection*{Turunan Fungsi}

Definisi turunan:
$$f'(x) = \lim_{h \to 0} \frac{f(x+h) - f(x)}{h}$$

\emtcmd{\$showev('limit(((x+h)\^{}n-x\^{}n)/h,h,0)) \textcolor{emtcomment}{// turunan x\^{}n}}

\textbf{Penjelasan Matematis (Teorema Binomial):} \\
Berdasarkan ekspansi binomial, $(x+h)^n = x^n + n x^{n-1} h + \frac{n(n-1)}{2} x^{n-2} h^2 + \dots + h^n$. Jika dikurangi $x^n$ lalu dibagi $h$, diperoleh: $n x^{n-1} + \frac{n(n-1)}{2} x^{n-2} h + \dots + h^{n-1}$. Saat $h \to 0$, seluruh suku yang masih memuat $h$ bernilai $0$, sehingga tersisa $n x^{n-1}$.

\emtcmd{\$showev('limit((sin(x+h)-sin(x))/h,h,0)) \textcolor{emtcomment}{// turunan sin(x)}}

\textbf{Penjelasan Matematis (Rumus Jumlah Dua Sudut):} \\
Menggunakan identitas $\sin(x+h) = \sin(x)\cos(h) + \cos(x)\sin(h)$, bentuk limit menjadi:
$$\lim_{h \to 0} \left[\sin(x)\frac{\cos(h)-1}{h} + \cos(x)\frac{\sin(h)}{h}\right]$$
Karena $\lim_{h \to 0} \frac{\cos(h)-1}{h} = 0$ dan $\lim_{h \to 0} \frac{\sin(h)}{h} = 1$, hasilnya adalah $\cos(x)$.

\emtcmd{\$showev('limit((log(x+h)-log(x))/h,h,0)) \textcolor{emtcomment}{// turunan log(x)}}

\textbf{Penjelasan Matematis (Sifat Logaritma):} \\
Menggunakan sifat $\log(x+h) - \log(x) = \log\left(\frac{x+h}{x}\right) = \log\left(1 + \frac{h}{x}\right)$, bentuk limit menjadi:
$$\lim_{h \to 0} \frac{1}{h}\log\left(1 + \frac{h}{x}\right) = \lim_{h \to 0} \log\left(\left(1 + \frac{h}{x}\right)^{1/h}\right)$$
Misalkan $u = \frac{h}{x}$, saat $h \to 0$ maka $u \to 0$ dan $\frac{1}{h} = \frac{1}{ux}$, sehingga diperoleh:
$$\frac{1}{x}\log\left(\lim_{u \to 0} (1 + u)^{1/u}\right) = \frac{1}{x}\log(e) = \frac{1}{x}$$

\emtcmd{\$showev('limit((1/(x+h)-1/x)/h,h,0)) \textcolor{emtcomment}{// turunan 1/x}}
\emtcmd{\$showev('limit(factor((E\^{}(x+h)-E\^{}x)/h),h,0)) \textcolor{emtcomment}{// turunan e\^{}x}}

\emtcmd{function f(x) \&= x\^{}x}
\emtcmd{\&assume(x>0); \$showev('limit((f(x+h)-f(x))/h,h,0)) \textcolor{emtcomment}{// turunan x\^{}x}}
\emtcmd{\&forget(x>0) \textcolor{emtcomment}{// menghapus asumsi agar kembali ke pengaturan awal}}

\emtcmd{function f(x) \&= sinh(x)}
\emtcmd{function df(x) \&= limit((f(x+h)-f(x))/h,h,0); \$df(x)}
\emtcmd{plot2d(["f(x)","df(x)"],-pi,pi,color=[blue,red]):}

\emtcmd{function f(x) \&= sin(3*x\^{}5+7)\^{}2}
\emtcmd{diff(f,3), diffc(f,3)}

\textbf{Perbedaan \texttt{diff} dan \texttt{diffc}:}
\begin{itemize}
    \item \texttt{diff(f, x)} menghitung turunan numerik menggunakan selisih maju/mundur standar.
    \item \texttt{diffc(f, x)} menghitung turunan numerik menggunakan selisih pusat (central difference), yakni $\frac{f(x+h) - f(x-h)}{2h}$, yang umumnya memiliki galat lebih kecil (lebih akurat).
\end{itemize}

\subsubsection*{Latihan Turunan 5 Jenis Fungsi}

\textbf{1. Fungsi Polinomial:} $p_1(x) = 2x^3 - 5x + 3$
\emtcmd{function p1(x) \&= 2*x\^{}3 - 5*x + 3}
\emtcmd{function dp1(x) \&= limit((p1(x+h) - p1(x))/h, h, 0); \$dp1(x)}
\emtcmd{plot2d(["p1(x)", "dp1(x)"], -2, 2, color=[blue, red], title="Polinomial \& Turunannya"):}

\textbf{2. Fungsi Rasional:} $p_2(x) = \frac{x - 1}{x + 2}$
\emtcmd{function p2(x) \&= (x - 1)/(x + 2)}
\emtcmd{function dp2(x) \&= ratsimp(limit((p2(x+h) - p2(x))/h, h, 0)); \$dp2(x)}
\emtcmd{plot2d(["p2(x)", "dp2(x)"], -1, 5, color=[blue, red], title="Rasional \& Turunannya"):}

\textbf{3. Fungsi Trigonometri:} $p_3(x) = \cos(2x) + \sin(x)$
\emtcmd{function p3(x) \&= cos(2*x) + sin(x)}
\emtcmd{function dp3(x) \&= limit((p3(x+h) - p3(x))/h, h, 0); \$dp3(x)}
\emtcmd{plot2d(["p3(x)", "dp3(x)"], -pi, pi, color=[blue, red], title="Trigonometri \& Turunannya"):}

\textbf{4. Fungsi Eksponensial:} $p_4(x) = e^{2x}$
\emtcmd{function p4(x) \&= exp(2*x)}
\emtcmd{function dp4(x) \&= limit(factor((p4(x+h) - p4(x))/h), h, 0); \$dp4(x)}
\emtcmd{plot2d(["p4(x)", "dp4(x)"], -1.5, 1.5, color=[blue, red], title="Eksponensial \& Turunannya"):}

\textbf{5. Fungsi Akar (Irasional):} $p_5(x) = \sqrt{x + 3}$
\emtcmd{function p5(x) \&= sqrt(x + 3)}
\emtcmd{\&assume(x > -3);}
\emtcmd{function dp5(x) \&= limit((p5(x+h) - p5(x))/h, h, 0); \$dp5(x)}
\emtcmd{\&forget(x > -3);}
\emtcmd{plot2d(["p5(x)", "dp5(x)"], -2.5, 5, color=[blue, red], title="Fungsi Akar \& Turunannya"):}

\subsection*{Integral dan Aplikasinya}

Integral tak tentu dan integral tentu menggunakan fungsi Maxima \texttt{"integrate"}.

\emtcmd{\$showev('integrate(1/(1+x\^{}2),x))}
\emtcmd{\$showev('integrate(sin(x),x,0,pi))}
\emtcmd{\$showev('integrate(x*exp(-x),x,0,1))}

\subsubsection*{Latihan 4A: Integral 5 Fungsi Integrabel \& 3 Fungsi Non-Integrabel Elementer}

\textbf{A. 5 Fungsi Integrabel Elementer:}

\textbf{1. Polinomial:} $\int_0^3 (3x^2 - 4x + 2) \, dx$
\emtcmd{\$showev('integrate(3*x\^{}2 - 4*x + 2, x))}
\emtcmd{\$showev('integrate(3*x\^{}2 - 4*x + 2, x, 0, 3))}
\emtcmd{plot2d("3*x\^{}2 - 4*x + 2", -0.5, 3.5); x\_ar=linspace(0,3,100); plot2d(x\_ar, 3*x\_ar\^{}2-4*x\_ar+2, >filled, style="/", fillcolor=green, >add):}

\textbf{2. Rasional:} $\int_0^1 \frac{4}{1 + x^2} \, dx$
\emtcmd{\$showev('integrate(4/(1 + x\^{}2), x))}
\emtcmd{\$showev('integrate(4/(1 + x\^{}2), x, 0, 1))}
\emtcmd{plot2d("4/(1 + x\^{}2)", -0.5, 1.5); x\_ar=linspace(0,1,100); plot2d(x\_ar, 4/(1+x\_ar\^{}2), >filled, style="/", fillcolor=cyan, >add):}

\textbf{3. Trigonometri:} $\int_0^{\pi/2} \sin(x)\cos(x) \, dx$
\emtcmd{\$showev('integrate(sin(x)*cos(x), x))}
\emtcmd{\$showev('integrate(sin(x)*cos(x), x, 0, pi/2))}
\emtcmd{plot2d("sin(x)*cos(x)", 0, pi); x\_ar=linspace(0,pi/2,100); plot2d(x\_ar, sin(x\_ar)*cos(x\_ar), >filled, style="/", fillcolor=yellow, >add):}

\textbf{4. Eksponensial:} $\int_0^2 2e^{-x} \, dx$
\emtcmd{\$showev('integrate(2*exp(-x), x))}
\emtcmd{\$showev('integrate(2*exp(-x), x, 0, 2))}
\emtcmd{plot2d("2*exp(-x)", -0.2, 2.5); x\_ar=linspace(0,2,100); plot2d(x\_ar, 2*exp(-x\_ar), >filled, style="/", fillcolor=orange, >add):}

\textbf{5. Logaritma (Integral Parsial):} $\int_1^e \log(x) \, dx$
\emtcmd{\$showev('integrate(log(x), x))}
\emtcmd{\$showev('integrate(log(x), x, 1, exp(1)))}
\emtcmd{plot2d("log(x)", 0.5, 3); x\_ar=linspace(1,exp(1),100); plot2d(x\_ar, log(x\_ar), >filled, style="/", fillcolor=magenta, >add):}

\medskip
\textbf{B. 3 Fungsi Non-Integrabel Elementer (Integrasi Numerik):}

\textbf{1. Fungsi Gauss:} $e^{-x^2}$ pada $[0, 2]$
\emtcmd{integrate("exp(-x\^{}2)", 0, 2)}
\emtcmd{plot2d("exp(-x\^{}2)", -0.5, 2.5); x\_ar=linspace(0,2,100); plot2d(x\_ar, exp(-x\_ar\^{}2), >filled, style="|", fillcolor=green, >add):}

\textbf{2. Fungsi Sinus Kardinal:} $\frac{\sin(x)}{x}$ pada $[0.001, \pi]$
\emtcmd{integrate("sin(x)/x", 0.001, pi)}
\emtcmd{plot2d("sin(x)/x", 0, 4); x\_ar=linspace(0.001,pi,100); plot2d(x\_ar, sin(x\_ar)/x\_ar, >filled, style="|", fillcolor=blue, >add):}

\textbf{3. Fungsi Fresnel Sinus:} $\sin(x^2)$ pada $[0, 2]$
\emtcmd{integrate("sin(x\^{}2)", 0, 2)}
\emtcmd{plot2d("sin(x\^{}2)", 0, 2.5); x\_ar=linspace(0,2,100); plot2d(x\_ar, sin(x\_ar\^{}2), >filled, style="|", fillcolor=red, >add):}

\subsubsection*{Latihan 4B: Luas Daerah Antara Dua Kurva, Volume Benda Putar, dan Panjang Kurva}

\textbf{1. Luas daerah antara $y = 4 - x^2$ dan $y = x + 2$}
\emtcmd{titik\_potong \&= solve(4 - x\^{}2 = x + 2, x); \$titik\_potong}
\emtcmd{\$showev('integrate((4 - x\^{}2) - (x + 2), x, -2, 1))}
\emtcmd{plot2d(["4 - x\^{}2", "x + 2"], -2.5, 2, color=[blue, red]);}
\emtcmd{x\_lu = linspace(-2, 1, 200); xi\_lu = flipx(x\_lu);}
\emtcmd{plot2d(x\_lu|xi\_lu, (4 - x\_lu\^{}2)|(xi\_lu + 2), >filled, style="/", fillcolor=cyan, >add):}

\textbf{2. Volume benda putar $y = \sqrt{x}$ diputar mengelilingi sumbu-$x$ pada $[0, 4]$} \\
Rumus: $V = \pi \int_0^4 (\sqrt{x})^2 \, dx$
\emtcmd{\$showev('integrate(pi*(sqrt(x))\^{}2, x, 0, 4))}
\emtcmd{plot3d("x\^{}2 + y\^{}2", xmin=0, xmax=4, ymin=-2, ymax=2, title="Visualisasi Paraboloida Benda Putar"):}

\textbf{3. Panjang kurva $y = \frac{2}{3}x^{3/2}$ dari $x = 0$ sampai $x = 3$}
\emtcmd{function kurva\_s(x) \&= (2/3)*x\^{}(3/2)}
\emtcmd{\$showev('integrate(sqrt(1 + diff(kurva\_s(x), x)\^{}2), x, 0, 3))}
\emtcmd{plot2d("kurva\_s(x)", 0, 3.5, title="Kurva y = (2/3)x\^{}(3/2) pada [0,3]"):}

\subsubsection*{Latihan 4C: Empat Representasi Kurva \& Kurvatur (Kelengkungan)}

Kurva pilihan: Kurva Katenari $y = \cosh(x)$

\textbf{a. Kartesius:} $y = \cosh(x)$
\emtcmd{function y\_kart(x) \&= cosh(x)}
\emtcmd{function k\_kart(x) \&= trigsimp(diff(y\_kart(x), x, 2)/(1 + diff(y\_kart(x), x)\^{}2)\^{}(3/2)); \$'k\_kart(x)=k\_kart(x)}

\textbf{c. Parametrik:} $x(t) = t$, $y(t) = \cosh(t)$
\emtcmd{function xp(t) \&= t; function yp(t) \&= cosh(t);}
\emtcmd{function k\_param(t) \&= trigsimp((diff(xp(t),t)*diff(yp(t),t,2) - diff(xp(t),t,2)*diff(yp(t),t))/(diff(xp(t),t)\^{}2 + diff(yp(t),t)\^{}2)\^{}(3/2)); \$'k\_param(t)=k\_param(t)}

\textbf{d. Implisit:} $F(x, y) = \cosh(x) - y = 0$
\emtcmd{function F\_imp(x, y) \&= cosh(x) - y;}
\emtcmd{Fx \&= diff(F\_imp(x,y), x); Fxx \&= diff(F\_imp(x,y), x, 2);}
\emtcmd{Fy \&= diff(F\_imp(x,y), y); Fyy \&= diff(F\_imp(x,y), y, 2); Fxy \&= diff(diff(F\_imp(x,y), x), y);}
\emtcmd{function k\_imp(x) \&= trigsimp((Fy\^{}2*Fxx - 2*Fx*Fy*Fxy + Fx\^{}2*Fyy)/(Fx\^{}2 + Fy\^{}2)\^{}(3/2)); \$'k\_imp(x)=k\_imp(x)}

\textbf{b. Koordinat Kutub (Spiral Logaritmik $r(\theta) = e^{\theta}$):}
\emtcmd{function r\_kut(th) \&= exp(th);}
\emtcmd{function k\_kut(th) \&= radcan((r\_kut(th)\^{}2 + 2*diff(r\_kut(th),th)\^{}2 - r\_kut(th)*diff(r\_kut(th),th,2))/(r\_kut(th)\^{}2 + diff(r\_kut(th),th)\^{}2)\^{}(3/2)); \$'k\_kut(th)=k\_kut(th)}

Plot Katenari $y = \cosh(x)$, kurva kurvatur $k(x)$, radius kelengkungan $1/k(x)$, dan lingkaran oskulasi di puncak $(0, 1)$: \\
Di $x = 0$: $y(0) = 1$, $k(0) = 1$, radius $R = 1$, sehingga pusat lingkaran oskulasi di $(0, 2)$.
\emtcmd{plot2d(["cosh(x)", "1/cosh(x)\^{}2"], -2, 2, color=[blue, red]);}
\emtcmd{plot2d("cosh(x)\^{}2", -1.2, 1.2, color=green, >add);}
\emtcmd{plot2d("cos(x)", "2 + sin(x)", xmin=0, xmax=2*pi, color=magenta, >add, title="Katenari, Kurvatur, Radius, \& Lingkaran Oskulasi"):}

\subsection*{Barisan dan Deret}

Contoh perhitungan barisan, iterasi, dan deret tak hingga dengan EMT dan Maxima.

\emtcmd{1:10 \textcolor{emtcomment}{// barisan sederhana}}
\emtcmd{q=1.05; iterate("x*q",1000,n=10)' \textcolor{emtcomment}{// pertumbuhan 5\% selama 10 periode}}
\emtcmd{sequence("x[n-1]+x[n-2]",[1,1],15) \textcolor{emtcomment}{// barisan Fibonacci}}
\emtcmd{iterate("cos(x)",1) \textcolor{emtcomment}{// iterasi titik tetap x = cos(x)}}
\emtcmd{iterate("(x+2/x)/2",2) \textcolor{emtcomment}{// metode Newton untuk akar(2)}}

Deret Taylor dan Jumlah Deret Tak Hingga:
\emtcmd{\$'sum(k,k,1,n) = factor(ev(sum(k,k,1,n),simpsum=true))}
\emtcmd{\$'sum(1/x\^{}2,x,1,inf) = ev(sum(1/x\^{}2,x,1,inf),simpsum=true)}
\emtcmd{\&assume(abs(x)<1); \$'sum(a*x\^{}k,k,0,inf) = ev(sum(a*x\^{}k,k,0,inf),simpsum=true); \&forget(abs(x)<1)}
\emtcmd{\$'e\^{}x = taylor(exp(x),x,0,10)}
\emtcmd{\$'log(x) = taylor(log(x),x,1,10)}

\end{document}